%%%%%%%%%%%%%%%%%%%%%%%%%%%%%%%%%%%%%%%%%%%%%%%%%%%%%%%%%%%%%%%%%%%%%%%%%%%%%%%%
%% Capitulo 7: Guia de Implantação, DevOps e Runbook Operacional
%% Sistema: PostoSmart - Automação e Gestão de Postos de Combustíveis
%% Autor: Diogo Santos Pires Jandiroba
%%%%%%%%%%%%%%%%%%%%%%%%%%%%%%%%%%%%%%%%%%%%%%%%%%%%%%%%%%%%%%%%%%%%%%%%%%%%%%%%

\chapter{Guia de Implantação, DevOps e Runbook Operacional}
\label{cap:implantacao_devops}

\section{Pipeline de Integração e Entrega Contínua (CI/CD)}

A esteira de automação garante que qualquer alteração seja testada, escaneada quanto a vulnerabilidades e empacotada em contêineres imutáveis antes da promoção para produção (\autoref{fig:pipeline_cicd_posto}).

% PLACEHOLDER STARUML v7.0: Pipeline CI/CD
% Ferramenta: StarUML v7.0 -> Model -> Add Diagram -> Activity Diagram
% Arquivo: Imagens/devops_pipeline_posto.png (Exportar em PNG 300 DPI ou PDF vetorial)
% Descomente a linha abaixo apos salvar o arquivo no diretorio Imagens/:
% \incluirdiagrama{Imagens/devops_pipeline_posto.png}{Pipeline Automatizada de CI/CD do PostoSmart}{fig:pipeline_cicd_posto}
\begin{figure}[htbp]
    \centering
    \fbox{\parbox{0.9\textwidth}{\centering\vspace{1cm}%
        \textbf{[Pipeline de Integração e Entrega Contínua (CI/CD) -- StarUML v7.0]}\\%
        \textit{Modelar no StarUML v7.0 e exportar para: \texttt{Imagens/devops_pipeline_posto.png}}\\%
        \small Menu: \texttt{File -> Export Diagram as -> PNG...} (Fundo branco, alta resolução 300 DPI)\\%
        \vspace{0.3cm}%
        \footnotesize Para ativar a imagem real, descomente no arquivo \texttt{.tex}:\\
        \texttt{\textbackslash incluirdiagrama\{Imagens/devops_pipeline_posto.png\}\{Pipeline Automatizada de CI/CD do PostoSmart\}\{fig:pipeline_cicd_posto\}}%
    \vspace{1cm}}}
    \caption{Pipeline Automatizada de CI/CD do PostoSmart (Placeholder StarUML v7.0)}
    \label{fig:pipeline_cicd_posto}
\end{figure}

\section{Configuração de Contêineres e Execução Local (Docker Compose)}

O appliance de borda (\textit{Edge Server}) opera com uma pilha mínima e isolada em contêineres:

\begin{lstlisting}[language=bash]
# docker-compose.yml - Servidor de Pista Edge Local
version: '3.8'
services:
  postosmart-core:
    image: ghcr.io/diogojandiroba1/postosmart-core:1.0.0
    restart: always
    network_mode: host
    volumes:
      - /opt/postosmart/data:/app/data
      - /dev/ttyUSB0:/dev/ttyUSB0 # Acesso direto ao conversor RS-485
    environment:
      - CONCENTRADOR_IP=192.168.1.50
      - CONCENTRADOR_PORT=2001
      - SQLITE_DB_PATH=/app/data/postosmart.db
      - SEFAZ_ENV=PRODUCAO

  postosmart-sync:
    image: ghcr.io/diogojandiroba1/postosmart-sync:1.0.0
    restart: always
    depends_on:
      - postosmart-core
    environment:
      - CLOUD_API_URL=https://api.postosmart.com.br/v1/sync
      - CLIENT_CERT_PATH=/etc/ssl/posto_client.crt
\end{lstlisting}

\section{Runbook Operacional e Procedimentos de Contingência}

\subsection{Procedimento 01: Falha de Comunicação com o Concentrador de Bombas}
\paragraph{Sintoma:} O PDV exibe alerta vermelho: \texttt{"CONCENTRADOR INACESSÍVEL -- PISTA OPERANDO EM MODO AUTÔNOMO"}.
\paragraph{Ação Imediata do Operador / Suporte:}
\begin{enumerate}
    \item Verificar o LED indicativo de \texttt{LINK / DATA} no painel frontal do concentrador físico na parede;
    \item Confirmar a integridade do cabo Ethernet UTP Cat6 que liga o concentrador ao switch local;
    \item No terminal do servidor de borda, executar: \texttt{curl -I http://192.168.1.50:2001};
    \item Caso não haja resposta após 1 minuto, reiniciar o serviço do driver:
    \begin{lstlisting}[language=bash]
    docker compose restart postosmart-core
    \end{lstlisting}
    \item Se a falha persistir, chavear a chave comutadora física do concentrador para o modo \textit{BYPASS MANUAL}, permitindo que os frentistas abasteçam anotando os encerrantes mecânicos nas bombas.
\end{enumerate}

\subsection{Procedimento 02: Reconciliação e Despejo de NFC-e em Contingência}
Após a restauração da conectividade com a internet, o serviço de sincronização dispara a rotina:
\begin{lstlisting}[language=bash]
# Disparo manual de transmissao em lote de notas fiscais represadas
docker compose exec postosmart-core postosmart-cli flush-contingencia
\end{lstlisting}
Todas as notas autorizadas em contingência offline (RN03) são transmitidas em lote para a SEFAZ, gerando os protocolos definitivos e atualizando o status na nuvem.
